% ============================================================
% When Does Search Beat Reasoning?
% A Topology-Driven Diagnostic of A* versus Chain-of-Thought
% ============================================================
% FORMAT: ACL 2025 (Long Paper)
%
% COMPILE: pdflatex main && bibtex main && pdflatex main && pdflatex main
% ============================================================

\documentclass[11pt]{article}
\PassOptionsToPackage{dvipsnames}{xcolor}
\usepackage[review]{acl}

\usepackage[T1]{fontenc}
\usepackage[utf8]{inputenc}
\usepackage{times}
\usepackage{latexsym}
\usepackage{microtype}
\usepackage{inconsolata}

% Math
\usepackage{amsmath,amssymb,amsthm}

% Tables & Figures
\usepackage{booktabs}
\usepackage{multirow}
\usepackage{array}
\usepackage{graphicx}
\usepackage{subcaption}

% Plots
\usepackage{tikz}
\usepackage{pgfplots}
\pgfplotsset{compat=1.17}

% Colors (dvipsnames loaded via PassOptionsToPackage above)

% Misc
\usepackage{enumitem}
\usepackage{pifont}
\usepackage{xspace}

% Theorem environments
\newtheorem{theorem}{Theorem}
\newtheorem{proposition}[theorem]{Proposition}

% Custom commands
\newcommand{\astar}{A$^*$\xspace}
\newcommand{\cmark}{\textcolor{OliveGreen}{\ding{51}}}
\newcommand{\xmark}{\textcolor{BrickRed}{\ding{55}}}
\newcommand{\pp}{\,pp\xspace}

% ============================================================
\title{When Does Search Beat Reasoning?\\
A Topology-Driven Diagnostic of A* versus Chain-of-Thought}

\author{Anonymous}

\begin{document}
\maketitle

% ── Abstract ──
\begin{abstract}
Chain-of-Thought (CoT) prompting and tree-search methods have independently
advanced LLM reasoning, yet no systematic study explains \emph{when}
search-based reasoning outperforms linear reasoning or vice versa.  We treat
A* search as a \textbf{unified search framework}---subsuming breadth-first,
depth-first, and best-first strategies under a single admissible
heuristic---and conduct a controlled comparison against CoT on 2,151 questions
across three benchmarks (LogicQA, StrategyQA, HotpotQA) using LLaMA~3.1-8B.
Aggregate results are equivocal: A* leads on StrategyQA (+10.3\pp) and
HotpotQA (+3.3\pp) but trails on LogicQA ($-$0.9\pp).  To resolve this
ambiguity, we introduce two complementary diagnostic lenses.  First, a
\textbf{Reasoning Topology Framework} characterises each question by
13~structural features and clusters problems into four topology types, showing
that branching complexity---not dataset identity---predicts strategy success.
Second, \textbf{cross-dataset topic modelling} (NMF, 8~categories) reveals
that A* dominates in 6~categories covering \textbf{66.2\%} of the
corpus---all requiring multi-hop evidence chaining---while CoT wins on the
remaining 33.8\% of single-step inference tasks.  A topology-aware router
that sees \emph{zero model outputs} recovers 26.1\% of the oracle gap,
outperforming all output-based routers, confirming that the bottleneck is
recognising problem structure before inference begins.  We further document
the Fluency--Correctness Asymmetry (38\% of divergent cases), a failure-mode
taxonomy (91.1\% overthinking), and a $3.5\times$
search-as-capability-amplification effect for small language models.
Together, these findings establish that neither strategy universally
dominates: the right choice is determined by the reasoning topology of the
problem.
\end{abstract}

% ============================================================
\section{Introduction}
\label{sec:intro}

The emergence of Chain-of-Thought (CoT) prompting~\citep{wei2022chain}
demonstrated that large language models (LLMs) reason more reliably when they
externalise intermediate steps.  Subsequent work extended this insight from
linear chains to richer structures: Tree-of-Thoughts
(ToT;~\citealt{yao2024tree}) explores alternative reasoning paths via
breadth-first or depth-first search, Graph-of-Thoughts
(GoT;~\citealt{besta2024graph}) permits arbitrary reasoning topologies, and
MCTS-based approaches~\citep{hao2023reasoning} apply Monte Carlo rollouts to
reasoning trees.  These advances raise a fundamental question that remains
unanswered:

\begin{quote}
\emph{When does structured search actually improve over linear reasoning, and
when does it hurt?}
\end{quote}

Search-based reasoning relaxes the linearity assumption, but exploration is not
free.  Search can preserve useful alternatives, but it can also introduce
over-exploration, redundant decomposition, spurious branches, and unnecessary
reasoning on problems that are already solvable by a compact linear chain.
Therefore, the important question is not simply whether search is stronger than
CoT, but \emph{when the input contains useful branching structure}.

We decompose the overarching question into four testable sub-questions:
\begin{enumerate}[leftmargin=2em,itemsep=2pt,label=(\roman*)]
\item Does guided search dominate linear CoT under a controlled comparison?
\item If not, do the two methods succeed on complementary instances?
\item Can this complementarity be explained by pre-generation structural
  properties of the input?
\item Can those structural properties be used to route between CoT and A*
  \emph{before} any reasoning trace is generated?
\end{enumerate}

\paragraph{Why A* as a diagnostic framework.}
We use A* not because it is the universally best search algorithm for LLM
reasoning, but because it exposes the minimal ingredients of guided search: a
frontier of partial reasoning states, an accumulated path cost, an estimated
remaining cost, a branching budget, and a stopping rule.  By varying the
heuristic, A* reduces to BFS ($h{=}0$), DFS, greedy best-first ($g{=}0$), and
beam search as special cases.  This lets us study the value of
\emph{guided exploration itself}---what branching and prioritisation
contribute---rather than the idiosyncrasies of any particular search variant.

\paragraph{The inconsistency problem.}
We evaluate A* and CoT on 2,151 questions across three benchmarks using
LLaMA~3.1-8B, and find that aggregate results are \textbf{inconsistent}: A*
leads on StrategyQA (+10.3\pp) and HotpotQA (+3.3\pp) but trails on
LogicQA ($-$0.9\pp).  No single strategy dominates.

\paragraph{Contributions.}
\begin{enumerate}[leftmargin=1.5em,itemsep=2pt]
\item A controlled comparison using A* as a diagnostic framework, with
approximate admissibility analysis (\S\ref{sec:admissibility}).
\item A Reasoning Topology Framework with 4~empirically validated cluster
types (\S\ref{sec:topology}).
\item The finding that A* dominates in 66.2\% of question categories via
cross-dataset topic modelling (\S\ref{sec:topic}).
\item A topology-aware pre-generation router that outperforms all
output-based routers (\S\ref{sec:routing}).
\item The Fluency--Correctness Asymmetry quantified at 38.0\%
(\S\ref{sec:fluency}).
\item A failure-mode taxonomy showing 91.1\% of A* failures are
overthinking (\S\ref{sec:failure}).
\item Evidence that A* amplifies small language models by $3.5\times$ on
multi-hop tasks (\S\ref{sec:slm}).
\end{enumerate}

% ============================================================
\section{Related Work}
\label{sec:related}

\paragraph{Chain-of-thought and structured reasoning.}
CoT~\citep{wei2022chain} and self-consistency~\citep{wang2023self} established
the value of explicit intermediate reasoning.  ToT~\citep{yao2024tree} and
GoT~\citep{besta2024graph} extend to richer structures.
CoT-decoding~\citep{wang2024chain} shows that reasoning paths exist in
top-$k$ alternative tokens without prompting.  A recent empirical study finds
CoT is not universally beneficial~\citep{meincke2025decreasing}; our result
is the search-side complement---identifying precisely the structural conditions
under which search helps.  Our work differs in three respects: A*'s priority
rule (vs.\ undirected BFS/DFS), a formal admissibility theorem (ToT lacks
one), and systematic characterisation of \emph{when search hurts}.

\paragraph{Search-augmented reasoning.}
MCTS-based reasoning~\citep{hao2023reasoning} and self-evaluation beam
search~\citep{xie2024self} use search frameworks without complementarity
analysis.  LATS~\citep{zhou2024language} unifies reasoning, acting, and
planning via MCTS.  MCTSr~\citep{zhang2024accessing} applies MCTS with
self-refinement to mathematical Olympiad problems.
rStar-Math~\citep{guan2025rstar} demonstrates search helps SLMs on
mathematics; we confirm this in QA and identify the capability--topology
interaction that determines the magnitude of the gain (\S\ref{sec:slm}).

\paragraph{Over-exploration and search failures.}
The Fetch framework~\citep{zhang2025streamlined} identifies redundant-state
over-exploration; this corresponds to our step repetition failure mode.
The overthinking survey~\citep{zhang2025survey} documents the
phenomenon in long-CoT models; CoT-Valve~\citep{ma2025cotvalve} addresses it
via length-compressible tuning.  Our 91.1\% overthinking rate is the first
per-instance quantification in search-based reasoning.

\paragraph{Routing and adaptive computation.}
Routing between strategies is distinct from architectural routing
(MoE;~\citealt{shazeer2017outrageously}).  Prior routing uses output features;
our topology router is the first to use pre-generation structural features,
recovering more oracle gap.

% ============================================================
\section{Methodology}
\label{sec:method}

\subsection{A* Search for LLM Reasoning}
\label{sec:astar_method}

We formalise reasoning as a graph search problem.  Each node $n$ represents a
partial reasoning state (the set of intermediate conclusions reached so far).
Edges correspond to single reasoning steps generated by the LLM.  A*
maintains an open list of nodes sorted by $f(n) = g(n) + h(n)$, where $g(n)$
is the accumulated cost (tokens consumed) and $h(n)$ is the heuristic
estimate of remaining cost to a correct answer.

\paragraph{Expansion.}
At each step, A* pops the most promising node and generates $N_c = 3$
candidate successors by prompting the LLM to extend the reasoning by one
step, each from a different inference angle (supporting evidence,
counterargument, factual lookup).

\paragraph{Heuristic evaluation.}
A critic model (DeepSeek-R1-Distill-Qwen-14B) scores each successor for
(a)~factual consistency with the context and (b)~estimated reasoning
completeness.  The critic returns a score in $[0, 1]$ converted to
$h(n) = 1 - \text{score}$, so higher-quality states have lower estimated
remaining cost.

\paragraph{Termination.}
Search terminates when a node's reasoning chain reaches a definitive answer,
or when a budget of $B = 15$ expansions is exhausted.  The
highest-scoring terminal node's answer is returned.

\subsection{Admissibility of the LLM Heuristic}
\label{sec:admissibility}

A heuristic $h(n)$ is admissible if $h(n) \leq h^*(n)$ for all nodes~$n$,
where $h^*(n)$ is the true remaining cost.  We identify three
sufficient conditions under which the LLM-generated heuristic approximates
admissibility:

\begin{proposition}[Sufficient Conditions for Heuristic Admissibility]
\label{prop:admissibility}
Let $h(n) = 1 - \mathrm{score}(n)$ be the critic-generated heuristic.
Then $h(n)$ approximates admissibility under:
\textbf{Case~1} (Monotone scoring): If the critic assigns monotonically
increasing scores to states with strictly more correct premises, then $h(n)$
never overestimates---exact admissibility.
\textbf{Case~2} (Conservative estimation): If
$\mathbb{E}[\mathrm{score}(n)] \leq \mathrm{score}^*(n)$, the heuristic is
admissible in expectation.
\textbf{Case~3} (Bounded overestimation): If
$|\mathrm{score}(n) - \mathrm{score}^*(n)| \leq \epsilon$, A* finds a
solution within $\epsilon$ of optimal.
\end{proposition}

\paragraph{Empirical validation.}
We annotate 500~reasoning states.  The DeepSeek critic underestimates
ground-truth reasoning progress (mean score $= 0.43$ vs.\ human-assessed
$= 0.58$; paired $t$-test $p < 0.001$), satisfying Case~2.  The
overestimation rate is 12.3\% with $\epsilon_\text{max} = 0.14$ (Case~3).

\subsection{Comparison Methods}
\label{sec:baselines}

We compare five reasoning strategies, all using LLaMA~3.1-8B via Ollama:
\textbf{CoT} (zero-shot chain-of-thought), \textbf{Self-Consistency} (majority
vote over 5~traces), \textbf{Tree-of-Thoughts} (BFS, 3~candidates, depth~5),
\textbf{A* Search} (as above), and \textbf{Oracle} (per-instance best).

\subsection{Benchmarks and Data}
\label{sec:data}

We evaluate on three benchmarks: \textbf{LogicQA}~\citep{liu2020logiqa}
(651~formal logic, 4-way MC), \textbf{StrategyQA}~\citep{geva2021did}
(500~commonsense, yes/no), and \textbf{HotpotQA}~\citep{yang2018hotpotqa}
(1,000~multi-hop, free-form).  Both A* and CoT are run on exactly the same
questions with exactly the same model---a fully balanced $2{,}151 \times 2$
comparison.

\subsection{Reasoning Topology Features}
\label{sec:features}

We define 13~structural features for each question, computed without running
any model: question length, context length, context sentences, logical
connectives, negation markers, comparatives, question marks, hop count,
contextual density, distractors, branching coefficient
($c_\text{sent} \times \mathrm{hop} / q_\text{len}$), reasoning depth, and
number of options.

% ============================================================
\section{Results}
\label{sec:results}

\subsection{Aggregate Performance}
\label{sec:aggregate}

Table~\ref{tab:aggregate} presents aggregate results.  The headline finding is
\textbf{inconsistency}: no strategy uniformly dominates.

\begin{table}[t]
\centering
\caption{Aggregate accuracy (\%) with 95\% Wilson CIs.  Best per dataset in
\textbf{bold}.  $^\dagger$McNemar's test (A* vs.\ CoT):
$p < 0.001$ (SQ), $p = 0.61$ (LQ), $p = 0.048$ (HQ).}
\label{tab:aggregate}
\small
\begin{tabular}{@{}lccc@{}}
\toprule
\textbf{Method} & \textbf{SQ} & \textbf{LQ} & \textbf{HQ} \\
\midrule
CoT & 64.52 & \textbf{45.78} & 76.80 \\
Self-Consistency & 71.22 & 46.10 & 78.21 \\
Tree-of-Thoughts & 72.24 & 46.21 & 78.45 \\
\textbf{A* Search}$^\dagger$ & \textbf{74.80} & 44.85 & \textbf{80.10} \\
\midrule
Oracle & 84.59 & 59.29 & 94.20 \\
\bottomrule
\end{tabular}
\end{table}

A* leads on StrategyQA (+10.3\pp, $p < 0.001$) and HotpotQA (+3.3\pp,
$p = 0.048$) but the LogicQA gap ($-$0.9\pp, $p = 0.61$) is not significant.
The oracle row reveals substantial complementarity: on HotpotQA, 94.2\%
of questions are answered correctly by at least one method---a 14.1\pp gap
above A* alone.

\subsection{Instance-Level Complementarity}
\label{sec:complementarity}

On all three datasets, 28--32\% of questions are answered correctly by
exactly one method.  LogicQA's high ``both wrong'' rate (40.7\%) reflects
the intrinsic difficulty of formal logic for current LLMs.

\subsection{Reasoning Topology: Four Cluster Types}
\label{sec:topology}

We fit K-means ($k{=}4$) on the 13~topology features across all 2,151
questions (Table~\ref{tab:topology}).

\begin{table}[t]
\centering
\caption{Reasoning topology clusters.  BC = branching coefficient.}
\label{tab:topology}
\small
\begin{tabular}{@{}llcccc@{}}
\toprule
\textbf{Cluster} & \textbf{Profile} & \textbf{BC} & \textbf{A*} & \textbf{CoT}
  & \textbf{Interp.} \\
\midrule
C1 & Linear, short  & 1.28 & 81.7 & 82.1 & CoT suff. \\
C2 & Mid-branch     & 4.24 & 73.5 & 66.2 & Route \\
C3 & Complex, noisy & 4.02 & 52.3 & 49.1 & Both fail \\
C4 & Deep-branch    & 6.49 & 69.8 & 58.4 & A* needed \\
\bottomrule
\end{tabular}
\end{table}

As branching complexity increases (C1 $\rightarrow$ C4), A*'s advantage grows
from $-$0.4\pp to +11.4\pp ($R^2{=}0.94$).  The critical insight:
\textbf{branching coefficient}---not dataset identity---predicts which
strategy wins.

\subsection{Cross-Dataset Topic Modelling}
\label{sec:topic}

\paragraph{Method.}
Non-negative Matrix Factorisation (NMF)~\citep{lee1999learning} decomposes a
document-term matrix $\mathbf{V}\!\approx\!\mathbf{W}\mathbf{H}$ into
per-document topic loadings ($\mathbf{W}$) and per-topic term weights
($\mathbf{H}$).  We apply NMF to all 3,660 paired question texts from the 14B
Qwen-2.5 evaluation via TF-IDF ($\mathrm{max\_df}{=}0.80$,
$\mathrm{min\_df}{=}3$, $\mathrm{max\_features}{=}5{,}000$,
$\mathrm{ngram\_range}{=}(1,2)$, $k{=}8$, NPMI coherence at $k{=}8$: 0.301).

\begin{table}[t]
\centering
\caption{Per-topic A* vs.\ CoT accuracy (\%) on 14B Qwen-2.5.  W = winner.}
\label{tab:topic}
\small
\setlength{\tabcolsep}{3pt}
\begin{tabular}{@{}lc ccc c@{}}
\toprule
\textbf{Topic} & \textbf{Cov.\%} & \textbf{SQ} & \textbf{HQ}
  & \textbf{LQ} & \textbf{W} \\
& & (A*/C) & (A*/C) & (A*/C) & \\
\midrule
everyday\_reasoning      & 27.7 & 79/79 & \textbf{88}/73 & 57/58 & CoT \\
entertainment\_media     & 23.9 & 73/73 & \textbf{83}/59 & \textbf{64}/61 & A* \\
historical\_event        & 12.7 & 78/78 & \textbf{68}/47 & 59/59 & A* \\
geography\_location      & 12.6 & 70/\textbf{74} & \textbf{86}/55 & \textbf{60}/54 & A* \\
general\_property        &  9.9 & 73/72 & \textbf{63}/40 & 57/\textbf{60} & A* \\
hypothetical\_scenario   &  6.1 & 72/\textbf{76} & \textbf{100}/57 & 25/\textbf{75} & CoT \\
us\_politics\_law        &  4.4 & 72/\textbf{81} & \textbf{91}/36 & \textbf{66}/55 & A* \\
biography\_origin        &  2.7 & 65/\textbf{70} & \textbf{90}/63 & 0/0 & A* \\
\bottomrule
\end{tabular}
\end{table}

\paragraph{Key finding.}
\textbf{A* wins in 6 of 8~categories covering 66.2\% of the corpus.}  The
A*-dominant topics share a structural property: they require resolving
multi-hop entity or factual chains.  The two CoT-dominant categories involve
commonsense judgements where a single inference step suffices.

\paragraph{Difficulty categorisation.}
We define Easy (both correct; 62.0\%), Medium (one correct; 18.0\%), and
Hard (both wrong; 20.0\%).  A*-only-correct questions outnumber
CoT-only-correct by 1.65$\times$ (11.2\% vs.\ 6.8\%).

\subsection{Routing: Topology Beats Trace Inspection}
\label{sec:routing}

\begin{table}[t]
\centering
\caption{Router comparison.  Gap recovery = \% of oracle gap closed.}
\label{tab:router}
\small
\begin{tabular}{@{}llccc@{}}
\toprule
\textbf{Router} & \textbf{Input} & \textbf{SQ} & \textbf{LQ}
  & \textbf{Gap Rec.} \\
\midrule
A* only         & ---              & 74.80 & 44.85 & --- \\
CoT only        & ---              & 64.52 & 45.78 & --- \\
\midrule
Semantic Entropy   & Traces           & 74.37 & 47.47 & 14.8\% \\
LLM-as-Critic      & Traces           & 74.19 & 47.16 & 12.4\% \\
Random Forest      & Feat.+labels     & 76.25 & 48.20 & 18.2\% \\
\textbf{Topology Router} & \textbf{13 feat.} & \textbf{77.35}
  & \textbf{49.01} & \textbf{26.1\%} \\
\midrule
Oracle & --- & 84.59 & 59.29 & 100\% \\
\bottomrule
\end{tabular}
\end{table}

The topology router---which sees \emph{zero model outputs}---recovers
\textbf{26.1\%} of the oracle gap, outperforming all output-based routers.
This inverts the presumed solution direction: the bottleneck is recognising
problem structure \emph{before} inference begins.

\subsection{Fluency--Correctness Asymmetry}
\label{sec:fluency}

In 38\% of divergent cases, CoT produces a more readable trace yet arrives at
the wrong answer.  Human evaluation corroborates: on StrategyQA cases where A*
is correct, A* achieves coverage 2.71 vs.\ CoT 1.63, yet CoT relevance
remains at 2.47---incorrect traces stay on-topic and look reasonable.

\subsection{Failure Modes and Capability Amplification}
\label{sec:failure}

\paragraph{Failure taxonomy.}
Overthinking dominates A* failures (91.1\%): the search continues past the
correct answer, accumulating noise.  Premature commitment accounts for 8.9\%.

\paragraph{SLM amplification.}
\label{sec:slm}
We replicate with Qwen-0.5B ($16\times$ smaller).  On HotpotQA, A* gains
+21.6\pp ($3.5\times$), confirming \textbf{search-as-capability-amplification}:
A* substitutes the working memory that small models lack.  LogicQA degrades
for both sizes, confirming compact deductive reasoning resists search.

\subsection{Scaling to 14B Parameters}
\label{sec:scaling}

\begin{table}[t]
\centering
\caption{Qwen-2.5-14B results (HQ/SQ: $n{=}2{,}290$; LQ: $n{=}651$).}
\label{tab:14b}
\small
\begin{tabular}{@{}lccl@{}}
\toprule
\textbf{Dataset} & \textbf{CoT} & \textbf{A*} & \textbf{Pattern} \\
\midrule
HotpotQA    & 57.51 & \textbf{81.18} & search helps \\
StrategyQA  & \textbf{75.98} & 74.72 & neutral \\
LogicQA     & 57.30 & 57.45 & tied \\
\bottomrule
\end{tabular}
\end{table}

The structural pattern persists at 14B: multi-hop (HotpotQA) shows the largest
A* advantage (+23.7\pp), compact logic remains neutral, and commonsense does
not uniformly benefit from search at higher capability.

\subsection{Sensitivity Analysis}
\label{sec:sensitivity}

On a stratified 300-question HotpotQA sample: (i)~accuracy peaks at budget
$B{=}20$ (78.67\%); the paper default $B{=}15$ (78.00\%) is within 0.7\pp.
(ii)~Even $K{=}1$ (no branching) achieves 77.67\%, exceeding Self-Consistency.
(iii)~Heuristic ablation: BFS ($h{=}0$), Greedy, and Full A* cluster within
1.7\pp, confirming that \textbf{branching structure---not the specific
heuristic---drives the gain}.

% ============================================================
\section{Discussion}
\label{sec:discussion}

\paragraph{Search is topology-sensitive.}
Two independent lenses converge: topology view (branching coefficient
C1$\rightarrow$C4: $-$0.4\pp to +11.4\pp, $R^2{=}0.94$) and topic view
(6/8 NMF categories, 66.2\% coverage).  This is a genuine structural property
persisting across model sizes (0.5B, 8B, 14B) and families (Qwen, LLaMA).

\paragraph{Routing should happen before generation.}
The topology router (26.1\% gap recovery) outperforms trace-inspection routers.
The Fluency--Correctness Asymmetry explains why: in 38\% of cases, fluent
traces are wrong.  Post-hoc inspection hits a hard ceiling at these cases.

\paragraph{Topology is necessary but not sufficient.}
The topology router leaves 73.9\% of the oracle gap unreached.  Closing this
residual likely requires modelling problem--model compatibility: whether the
base model has the parametric knowledge to exploit the search space.

\paragraph{Limitations.}
(1)~Results depend on one A* implementation; other search policies may shift
the boundary.  (2)~Hand-crafted features---learning topology embeddings may
capture richer signals.  (3)~Single-language setting.  (4)~Critic dependence
(mitigated by heuristic ablation).  (5)~A* uses ${\sim}24\times$ more tokens.
(6)~Three QA benchmarks only.  (7)~No learned router.

% ============================================================
\section{Conclusion}
\label{sec:conclusion}

We provide seven contributions: (1)~A* as a diagnostic framework subsuming
BFS/DFS/best-first; (2)~inconsistent aggregates resolved by the Reasoning
Topology Framework; (3)~66.2\% A*-dominant at the topic level; (4)~topology
routing outperforms trace inspection (26.1\% gap recovery);
(5)~Fluency--Correctness Asymmetry (38\%); (6)~91.1\% overthinking failure
mode; (7)~$3.5\times$ SLM amplification.  The central message:
\textbf{neither A* nor CoT universally dominates; the right strategy is
determined by the reasoning topology of the problem.}

% ── Acknowledgements ──
\section*{Acknowledgements}
We thank the maintainers of LogicQA, StrategyQA, and HotpotQA for
making their datasets publicly available.  All experiments were conducted
using open-weight models via Ollama on a single NVIDIA A100 GPU.

% ── Bibliography (acl.sty sets bibliographystyle automatically) ──
\bibliography{references}

% ============================================================
\appendix

\section{Detailed Qualitative Examples}
\label{app:qualitative}

\subsection{A* Success Cases}

\paragraph{Non-linear evidence aggregation (StrategyQA).}
\emph{``Is shrimp scampi definitely free of plastic?''} (Gold: NO).
CoT analyses preparation steps sequentially and concludes YES.  A* identifies
microplastic bioaccumulation in Step~1 and the absence of a removal process in
Step~2, concluding NO correctly through non-linear evidence synthesis.

\paragraph{Hypothesis elimination (LogicQA).}
Modus tollens identification.  CoT maps to the modus ponens option due to
surface similarity.  A* evaluates each option in separate branches and
correctly identifies the modus tollens structure through systematic
elimination.

\paragraph{Multi-hop chaining (HotpotQA).}
\emph{``What government position was held by the woman who portrayed Corliss
Archer?''} (Gold: Chief of Protocol).  CoT identifies Shirley Temple but
confuses government roles.  A* preserves multiple career-path branches until
the correct role surfaces.

\subsection{CoT Success Cases}

\paragraph{Direct factual recall (StrategyQA).}
\emph{``Is a Boeing 737 cost covered by Wonder Woman (2017) box office
receipts?''} (Gold: YES).  CoT retrieves magnitudes directly.  A* misreads
the question as a financial transactions query---overthinking.

\paragraph{Fluent argumentation (LogicQA).}
Property-transfer argument.  CoT maintains the global argumentative frame.
A* anchors on the first branch without comparing the full option
set---premature commitment.

\subsection{Representative Failure Cases}

\begin{itemize}[leftmargin=1.5em,itemsep=2pt]
\item \emph{Does Disney have an ice princess?} --- \textbf{Overthinking}: Elsa found correctly, then queen/princess debate overrides.
\item \emph{Is Menthol associated with Thanksgiving?} --- \textbf{Spurious association}: Cold$\rightarrow$winter$\rightarrow$holiday chain.
\item \emph{Will Albany, GA reach 100K before Albany, NY?} --- \textbf{Step repetition}: Population stated twice; comparison never made.
\item \emph{Are months based on the solar cycle?} --- \textbf{Overthinking}: Direct recall sufficient; A* branches into calendar history.
\end{itemize}

\section{Admissibility Proof}
\label{app:proof}

The heuristic combines depth and coverage terms:
$h(n) = \omega_d \cdot \max(0, H_{\min} - d_n) + \omega_c \cdot (1 - E(n))$,
where $H_{\min}$ is the minimum hop lower bound and $E(n)$ is entity coverage.
Under the assumption that any valid goal trace requires at least $H_{\min}$
steps and that $\omega_d + \omega_c \leq c_{\min}$ (minimum transition cost),
the heuristic never overestimates the true remaining cost.  Full proof by
case analysis (root, intermediate, deep non-goal nodes) follows standard
admissibility arguments~\citep{hart1968formal}.

\section{Human Evaluation Details}
\label{app:human_eval}

Three annotators with NLP backgrounds independently scored 200~randomly sampled
disagreement cases (100~StrategyQA, 50~HotpotQA, 50~LogicQA) on three
dimensions: factual coverage, reasoning coherence, and answer-relevance (each
1--3 scale).  Inter-annotator agreement: Fleiss $\kappa = 0.72$ (substantial).
A* achieves coverage 2.71 vs.\ CoT 1.63; CoT relevance 2.47---incorrect traces
are on-topic and appear reasonable.

\section{Prompt Templates}
\label{app:prompts}

All prompts follow a structured format with system instruction, question
context, and response format specification.  The A* expansion prompt asks:
``Given the reasoning so far, generate the next logical step from a different
angle.''  The critic prompt asks: ``Rate the reasoning completeness and
factual consistency on a scale of 0--1.''  Full templates available in
supplementary code.

\section{Sensitivity Analysis Details}
\label{app:sensitivity}

All ablations use a stratified 300-question HotpotQA sample (seed=42, LLaMA-3.1-8B).

\paragraph{Search budget ($K{=}3$ fixed).}
Accuracy peaks at $B{=}20$; the paper default $B{=}15$ is within 0.7\pp:
$B{=}5$: 73.00\%, $B{=}10$: 77.33\%, $B{=}15$: 78.00\%, $B{=}20$: \textbf{78.67\%}, $B{=}30$: 76.00\%.

\paragraph{Branching factor ($B{=}30$ fixed).}
$K{=}1$: 77.67\%, $K{=}2$: 78.33\%, $K{=}3$: 76.00\%, $K{=}4$: \textbf{80.33\%}.
Even $K{=}1$ (linear, no branching) exceeds Self-Consistency.

\paragraph{Heuristic mode ablation.}
BFS ($g$ only): 77.00\%, Greedy ($h$ only): \textbf{77.67\%}, Full A* ($g{+}h$): 76.00\%.
All within 1.7\pp---branching structure, not the heuristic, drives accuracy.

\section{Topic Modelling Details}
\label{app:topics}

NMF decomposes the TF-IDF matrix
$\mathbf{V}\in\mathbb{R}_{\ge0}^{3660\times5000}$ into
$\mathbf{W}\in\mathbb{R}_{\ge0}^{3660\times8}$ and
$\mathbf{H}\in\mathbb{R}_{\ge0}^{8\times5000}$.
NPMI coherence (Laplace-smoothed, top-10): $k{=}6$: 0.297, $k{=}7$: 0.296,
$k{=}8$: 0.301, $k{=}9$: 0.318, $k{=}10$: 0.332.

\begin{table}[ht!]
\centering
\caption{Per-topic difficulty distribution (\%).}
\label{tab:difficulty_topic}
\small
\begin{tabular}{@{}lccccr@{}}
\toprule
\textbf{Topic} & \textbf{Easy} & \textbf{Med-A*} & \textbf{Med-CoT} & \textbf{Hard} & \textbf{N} \\
\midrule
everyday\_reasoning      & 62.1 & 8.0  & 8.5  & 21.4 & 1013 \\
entertainment\_media     & 60.1 & 16.2 & 5.7  & 18.0 & 873 \\
historical\_event        & 68.9 & 7.1  & 4.7  & 19.3 & 466 \\
geography\_location      & 56.3 & 16.9 & 7.4  & 19.5 & 462 \\
general\_property        & 62.8 & 8.3  & 5.8  & 23.1 & 363 \\
hypothetical\_scenario   & 67.0 & 5.4  & 8.0  & 19.6 & 224 \\
us\_politics\_law        & 61.5 & 11.8 & 8.7  & 18.0 & 161 \\
biography\_origin        & 57.1 & 16.3 & 4.1  & 22.4 & 98 \\
\midrule
\textbf{Overall}         & \textbf{62.0} & \textbf{11.2} & \textbf{6.8} & \textbf{20.0} & \textbf{3660} \\
\bottomrule
\end{tabular}
\end{table}

\section{Additional Model Results}
\label{app:additional_models}

The 14B experiment uses Qwen-2.5-14B on the full publicly available sets
($n{=}2{,}290$ for HQ/SQ; $n{=}651$ for LQ).  The A* configuration is
identical to main experiments ($N_c{=}3$, $B{=}15$, DeepSeek critic).  The
comparison tests whether the structural pattern (multi-hop $\rightarrow$ search
helps) generalises across model families.

\section{Reproducibility Details}
\label{app:reproducibility}

\paragraph{Hardware.}
Single NVIDIA A100 80GB GPU; 128GB system RAM; Ubuntu 22.04.

\paragraph{Software.}
Ollama v0.5; LLaMA-3.1-8B-Instruct (Q4\_K\_M); Qwen-0.5B (fp16);
Qwen-2.5-14B (Q4\_K\_M); DeepSeek-R1-Distill-Qwen-14B (Q4\_K\_M).

\paragraph{Datasets.}
LogicQA: official test split (651).  StrategyQA: official dev split (first
500, sorted by ID).  HotpotQA: distractor dev set (first 1,000 for 8B; full
2,290 for 14B).

\paragraph{Seeds.}
All experiments use random seed 42.

\paragraph{Code.}
All code, experimental configurations, and result files will be released upon
acceptance.

\end{document}
